\documentclass[a4paper,11pt]{article}

% ---------- Préambule minimal ----------
\usepackage[utf8]{inputenc}
\usepackage[T1]{fontenc}
\usepackage[french]{babel}
\usepackage{amsmath,amssymb}
\usepackage{enumitem}
\usepackage{tikz}
\usepackage{tasks}
\usetikzlibrary{arrows.meta}
\usepackage[margin=2.5cm]{geometry}
\usepackage{needspace}
\usepackage{currfile}
\usepackage{fancyhdr}
\usepackage{lastpage}
\settasks{label=\arabic*.}

\pagestyle{empty}
\graphicspath{ {../Images/} }
\input{\currfiledir ../def.tex}
\def \Document {Exercices}

% Environnement "exercice" : \begin{exercice}{<numéro>} ... \end{exercice}
\newenvironment{exercice}[1]{%
  \Needspace{6\baselineskip}\par\bigskip\noindent\textbf{Exercice #1}\par\smallskip\noindent\ignorespaces
}{%
  \par
}

% ---------- Droites graduées ----------
% Espace vertical pour la rédaction (exercice 7)
\newcommand{\espace}{\rule[-2cm]{0pt}{0pt}}

% Exercice 1 : droite de -2 à 3, graduée tous les 0,25
\newcommand{\droiteun}[1]{%
  \begin{tikzpicture}[x=2.6cm,baseline=-0.5ex]
    \draw[{Latex[length=2mm]}-{Latex[length=2mm]}] (-2.3,0) -- (3.3,0);
    \foreach \i in {-2,...,3} {
      \draw (\i,0.12) -- (\i,-0.12);
      \node[below] at (\i,-0.12) {\small$\i$};
    }
    \foreach \x in {-1.75,-1.5,...,2.75} \draw (\x,0.06) -- (\x,-0.06);
    #1
  \end{tikzpicture}%
}

% Exercice 4 : de -5 à 5
\newcommand{\droitetrois}{%
  \begin{tikzpicture}[x=1.1cm,baseline=-0.5ex]
    \draw[-{Latex[length=2mm]}] (-5.7,0) -- (5.7,0);
    \foreach \i in {-5,...,5} {
      \draw (\i,0.08) -- (\i,-0.08);
      \node[below] at (\i,-0.08) {\scriptsize$\i$};
    }
  \end{tikzpicture}%
}

% Exercice 6 : de -10 à 10
\newcommand{\droitecinq}{%
  \begin{tikzpicture}[x=0.7cm,baseline=-0.5ex]
    \draw[{Latex[length=2mm]}-{Latex[length=2mm]}] (-10.6,0) -- (10.6,0);
    \foreach \i in {-10,...,10} {
      \draw (\i,0.08) -- (\i,-0.08);
      \node[below] at (\i,-0.08) {\scriptsize$\i$};
    }
  \end{tikzpicture}%
}

% Exercice 11 : de -2 à 10
\newcommand{\droiteonze}{%
  \begin{tikzpicture}[x=0.55cm,baseline=-0.5ex]
    \draw[{Latex[length=2mm]}-{Latex[length=2mm]}] (-2.6,0) -- (10.6,0);
    \foreach \i in {-2,...,10} {
      \draw (\i,0.08) -- (\i,-0.08);
      \$\bullet$ \quad Calculatrice autorisée \quaenode[below] at (\i,-0.08) {\scriptsize$\i$};
    }
  \end{tikzpicture}%
}

% Exercice 6 : un cas = énoncé + droite + zone de réponse
\newcommand{\casecinq}[2]{%
  \par\medskip\noindent
  \begin{minipage}{\linewidth}
    #1. $#2$\par\smallskip
    \begin{center}
      \droitecinq\\[1.2em]
      $I\cap J=\underline{\hspace{3.2cm}}\qquad\qquad I\cup J=\underline{\hspace{3.2cm}}$
    \end{center}
  \end{minipage}\par\bigskip
}

% Exercice 8 : chaîne d'équivalences à compléter
\newcommand{\chainevide}[1]{%
  En effet : $\begin{array}{rl}
    x\in #1 & \Leftrightarrow \underline{\hspace{3cm}}\\
    & \Leftrightarrow \underline{\hspace{3cm}}\\
    & \Leftrightarrow \underline{\hspace{3cm}}\\
    & \Leftrightarrow \underline{\hspace{3cm}}
  \end{array}$%
}

\begin{document}
\pagestyle{fancy}

\fancyhf{} % clear existing header/footer entries
\fancyhead[L]{ {\Niveau}}
\fancyhead[C]{ {\Chapitre}}
\fancyhead[R]{ {\Document}}
\fancyfoot[L]{ A. Fernandez }
\fancyfoot[R]{{\thepage\  / \pageref{LastPage}}}
%{
%  \normalsize
%  \textbf{\Niveau} \hfill \textbf{\Chapitre} \hfill \textbf{\Document}
%}

% =====================================================
\begin{exercice}{1}
\begin{enumerate}[label=\arabic*.]
\item Lire l'abscisse de chacun des points $A$, $B$, $C$, $D$ et $E$ placés sur la droite graduée ci-dessous.

\medskip
\begin{center}
\droiteun{\foreach \x/\n in {-1.75/A,-0.5/B,0.25/C,1.5/D,2.75/E} {
  \fill (\x,0) circle[radius=2pt];
  \node[above=3pt] at (\x,0) {$\n$};
}}
\end{center}

\begin{tasks}(3)
\task $A(\dots\dots)$
\task $B(\dots\dots)$
\task $C(\dots\dots)$
\task $D(\dots\dots)$
\task $E(\dots\dots)$
\end{tasks}

\item Placer le plus précisément possible, sur la droite graduée ci-dessous, les points $F$, $G$, $H$ et $K$ d'abscisses respectives :
\[
\frac{3}{4} \qquad -\frac{5}{4} \qquad 2{,}5 \qquad \frac{7}{3}
\]

\begin{center}
\droiteun{}
\end{center}
\end{enumerate}
\end{exercice}

% =====================================================
\begin{exercice}{2}
Traduire les inégalités suivantes par des appartenances à des intervalles.
\begin{tasks}(2)
\task $-2\leq x\leq 6$
\task $4<x<9$
\task $-1\leq x<10$
\task $12<x\leq 17{,}5$
\task $x\leq 7$
\task $8<x$
\task $x<-2$
\task $-1\leq x$
\end{tasks}
\end{exercice}

% =====================================================
\begin{exercice}{3}
Traduire les appartenances suivantes par des inégalités.
\begin{tasks}(2)
\task $x\in[-7\,;\,3]$
\task $x\in\,]0\,;\,1[$
\task $x\in\,]2\,;\,5]$
\task $x\in\,]-\infty\,;\,3[$
\task $x\in[-4\,;\,+\infty[$
\task $x\in[5\,;\,10[$
\task $x\in\,]-\infty\,;\,-1]$
\task $x\in\,]-2{,}5\,;\,0[$
\end{tasks}
\end{exercice}

% =====================================================
\begin{exercice}{4}
Représenter chacun des intervalles suivants sur une droite graduée.

\medskip
\noindent
\begin{tabular}{@{}p{2.6cm}l@{}}
1. $[-2\,;\,5]$ & \droitetrois \\[1.8em]
2. $]1\,;\,4]$ & \droitetrois \\[1.8em]
3. $[-1\,;\,+\infty[$ & \droitetrois \\[1.8em]
4. $]-\infty\,;\,3[$ & \droitetrois
\end{tabular}
\end{exercice}

% =====================================================
\begin{exercice}{5}
Compléter les pointillés par les symboles $\in$ ou $\notin$.
\begin{tasks}(2)
\task $3\ \dots\ [-4\,;\,6[$
\task $-2\ \dots\ [-1\,;\,4[$
\task $0\ \dots\ ]-2\,;\,1[$
\task $5\ \dots\ ]5\,;\,8]$
\task $7\ \dots\ ]6\,;\,7]$
\task $0\ \dots\ ]-\infty\,;\,0[$
\task $10^{-2}\ \dots\ ]0\,;\,+\infty[$
\task $10^{-5}\ \dots\ ]-\infty\,;\,0]$
\task $\pi\ \dots\ [3{,}1\,;\,3{,}2]$
\task $\dfrac{3}{7}\ \dots\ [0{,}5\,;\,2]$
\end{tasks}
\end{exercice}

% =====================================================
\begin{exercice}{6}
Dans chaque cas, représenter les intervalles $I$ et $J$ proposés sur la droite de deux couleurs différentes, puis déterminer leur intersection et leur union.

\casecinq{1}{I=[1\,;\,4]\ ;\ J=[3\,;\,7]}
\casecinq{2}{I=[-6\,;\,-1]\ ;\ J=[2\,;\,7{,}5]}
\casecinq{3}{I=\,]-\infty\,;\,0]\ ;\ J=[-3\,;\,5[}
\casecinq{4}{I=\,]-\infty\,;\,2[\ ;\ J=[1\,;\,+\infty[}
\casecinq{5}{I=\,]-\infty\,;\,4]\ ;\ J=[5\,;\,7[}
\casecinq{6}{I=[0\,;\,8[\ ;\ J=\,]0\,;\,3]}
\casecinq{7}{I=\,]-6\,;\,2[\ ;\ J=[2\,;\,+\infty[}
\end{exercice}

% =====================================================
\begin{exercice}{7}
Résoudre les inéquations et donner les solutions sous forme d'intervalles.
\begin{tasks}(3)
\task $x+3<12$\espace
\task $7x-4\leq 10$\espace
\task $6x+2\geq 56$\espace
\task $3x+6\leq 24$\espace
\task $-5x-11\geq -6$\espace
\task $-3x+2\leq 17$\espace
\end{tasks}
\end{exercice}

% =====================================================
\begin{exercice}{8}
Répondre à chaque fois en suivant l'exemple :

\medskip
\noindent
\begin{minipage}[t]{0.48\textwidth}
1. $x\in[3\,;\,5]\Leftrightarrow 2x+1\in[7\,;\,11]$

\smallskip
En effet : $\begin{array}{rl}
  x\in[3\,;\,5] & \Leftrightarrow 3\leq x\leq 5\\
  & \Leftrightarrow 6\leq 2x\leq 10\\
  & \Leftrightarrow 7\leq 2x+1\leq 11\\
  & \Leftrightarrow 2x+1\in[7\,;\,11]
\end{array}$
\end{minipage}
\hfill
\begin{minipage}[t]{0.48\textwidth}
2. $x\in\,]4\,;\,5]\Leftrightarrow 3x-2\in\underline{\hspace{2cm}}$

\smallskip
\chainevide{]4\,;\,5]}
\end{minipage}

\bigskip
\noindent
\begin{minipage}[t]{0.48\textwidth}
3. $x\in[1\,;\,2[\ \Leftrightarrow -3x+1\in\underline{\hspace{2cm}}$

\smallskip
\chainevide{[1\,;\,2[}
\end{minipage}
\hfill
\begin{minipage}[t]{0.48\textwidth}
4. $x\in[-5\,;\,2]\Leftrightarrow -2x-4\in\underline{\hspace{2cm}}$

\smallskip
\chainevide{[-5\,;\,2]}
\end{minipage}
\end{exercice}

% % =====================================================
% \begin{exercice}{9}
% Déterminer les valeurs suivantes :
% \begin{tasks}(3)
% \task $|-12|=\underline{\hspace{1.2cm}}$
% \task $|-4{,}5|=\underline{\hspace{1.2cm}}$
% \task $|7{,}2|=\underline{\hspace{1.2cm}}$
% \task $|4-10|=\underline{\hspace{1.2cm}}$
% \task $|-5-4|=\underline{\hspace{1.2cm}}$
% \task $|-7+9|=\underline{\hspace{1.2cm}}$
% \task $|10-(-6)|=\underline{\hspace{1.2cm}}$
% \task $\left|-\sqrt{2}\right|=\underline{\hspace{1.2cm}}$
% \task $\left|1-\sqrt{2}\right|=\underline{\hspace{1.2cm}}$
% \end{tasks}
% \end{exercice}
% 
% % =====================================================
% \begin{exercice}{10}
% On a placé quelques points sur la droite réelle ci-dessous :
% 
% \begin{center}
% \begin{tikzpicture}[x=2cm]
%   \draw[thick] (-1.8,0) -- (5,0);
%   % graduations tous les 0,5
%   \foreach \i in {-1.5,-1,...,4.5} \draw (\i,0.1) -- (\i,-0.1);
%   % origine et unité en rouge
%   \foreach \i in {0,1} {
%     \draw[red,thick] (\i,0.15) -- (\i,-0.15);
%     \node[red,below=3pt] at (\i,-0.1) {$\i$};
%   }
%   % points
%   \foreach \x/\n in {-1.5/K,-1/L,0.5/M,2.25/N,3.5/O,4.5/P} {
%     \draw[blue,thick] (\x,0.2) -- (\x,-0.2);
%     \node[blue,above=3pt] at (\x,0.2) {$\n$};
%   }
% \end{tikzpicture}
% \end{center}

% \begin{enumerate}[label=\arabic*.]
% \item Écrire en dessous de chaque point son abscisse.
% \item Déterminer les distances suivantes, en suivant l'exemple a.
%   \begin{tasks}[label=\alph*.](2)
%   \task LM $=|0{,}5-(-1)|=|1{,}5|=1{,}5$
%   \task KL $=$
%   \task KO $=$
%   \task OP $=$
%   \task ON $=$
%   \task PL $=$
%   \end{tasks}
% \end{enumerate}
% \end{exercice}

% =====================================================
% \begin{exercice}{9}
% Compléter les expressions suivantes, en vous aidant d'un schéma.
% 
% \medskip
% \noindent
% \begin{tabular}{@{}p{7.2cm}l@{}}
% 1. $x\in[3\,;\,7]\Leftrightarrow |x-5|\leq\dots$ & \droiteonze \\[1.6em]
% 2. $x\in[-1\,;\,5]\Leftrightarrow |x-\dots|\leq\dots$ & \droiteonze \\[1.6em]
% 3. $x\in[\dots\,;\,\dots]\Leftrightarrow |x-6|\leq 3$ & \droiteonze \\[1.6em]
% 4. $x\in[\dots\,;\,\dots]\Leftrightarrow |x-(-1)|\leq 1$ & \droiteonze \\[1.6em]
% 5. $x\in[3\,;\,\dots]\Leftrightarrow |x-5|\leq\dots$ & \droiteonze \\[1.6em]
% 6. $x\in[\dots\,;\,\dots]\Leftrightarrow |x|\leq 1$ & \droiteonze
% \end{tabular}
% \end{exercice}

% =====================================================
\begin{exercice}{9}
Pour chacun des nombres suivants, donner le plus petit sous-ensemble de $\mathbb{R}$ auquel il appartient parmi les ensembles du cours : $\mathbb{N}$, $\mathbb{Z}$, $\mathbb{D}$, $\mathbb{Q}$ et $\mathbb{R}$.
\begin{tasks}(4)
\task $15$
\task $\dfrac{2}{3}$
\task $4{,}8$
\task $-17$
\task $\dfrac{3}{2}$
\task $-4{,}275$
\task $\dfrac{12}{3}$
\task $\dfrac{-10}{5}$
\task $\sqrt{2}$
\task $\sqrt{9}$
\task $\sqrt{2}\times\sqrt{2}$
\task $\sqrt{2}-\sqrt{2}$
\end{tasks}
\end{exercice}

% =====================================================
\begin{exercice}{10}
Pour chacun des nombres suivants, donner un encadrement à $10^{-3}$ près, puis en déduire son arrondi à $10^{-3}$ près.
\begin{tasks}(3)
\task $\dfrac{1}{3}$
\task $0{,}7586$
\task $2{,}356\times 10^{-3}$
\end{tasks}
\end{exercice}

% =====================================================
\begin{exercice}{11}
Pour chacun des nombres suivants, donner un encadrement d'amplitude $10^{-2}$, puis en déduire son arrondi au centième.
\begin{tasks}(3)
\task $-327{,}426$
\task $\dfrac{5}{8}$
\task $4{,}526\,86\times 10^{2}$
\end{tasks}
\end{exercice}

% =====================================================
\begin{exercice}{12}
Vrai ou faux ? Si la réponse est \og faux \fg{}, donner un contre-exemple.
\begin{enumerate}[label=\arabic*.]
\item La différence de deux nombres entiers naturels est un nombre entier naturel.
\item Le produit de deux nombres entiers relatifs est un nombre entier naturel.
\item Le quotient d'un nombre entier relatif par un nombre entier relatif non nul est un nombre entier naturel.
\item Le quotient de deux nombres décimaux est un nombre décimal.
\item Le produit d'un nombre rationnel par un nombre entier relatif est un nombre rationnel.
\item La somme de deux nombres rationnels est un nombre rationnel.
\item Le quotient d'un nombre réel par un nombre réel non nul est un nombre rationnel.
\item Le produit de deux nombres irrationnels est un nombre irrationnel.
\end{enumerate}
\end{exercice}

\end{document}